\documentclass[11pt,a4paper]{article}
\usepackage[margin=24mm]{geometry}
\usepackage{amsmath,amssymb,amsthm,booktabs,hyperref}
\hypersetup{hidelinks,pdftitle={How fast can two replacement lonely runners be?}}
\newtheorem{theorem}{Theorem}
\newtheorem{lemma}[theorem]{Lemma}
\newcommand{\LR}{\operatorname{LR}}
\newcommand{\dist}[1]{\lVert #1\rVert}
\newcommand{\bad}{\mathcal B}
\newcommand{\good}{\mathcal U}
\setlength{\parskip}{2pt}
\title{\bfseries How fast can two replacement lonely runners be?}
\author{}
\date{\small Research draft \quad 26 November 2026}
\begin{document}
\maketitle
\begin{abstract}
Replace two distinct entries of $1,\ldots,n-1$ by integers $n\le u<w$.
If the resulting configuration is tight, we prove
$u=O(n\log\log n)$ and $w=O(n\log n)$. The order of the first
bound is attained along an infinite sequence. The second bound answers
the two-entry case of Zhang's Question 8.2. Simultaneous gaps in coprime
residues force a common primorial divisor. Testing several rational times
for each deletion then forces divisibility if $w$ is too large.
An exact search with no
prescribed limit on the inserted speeds finds three tight configurations
for $8\le n\le300$, all previously exhibited or predicted.
\end{abstract}

\section{The problem and the bounds}
On a circle of circumference one let $\dist{x}$ denote the distance to
the nearest integer. For a finite set of positive speeds $V$, put
$\LR(V)=\max_{t\in[0,1]}\min_{v\in V}\dist{vt}$. The lonely runner
conjecture predicts $\LR(V)\ge1/n$ for $n-1$ distinct positive integer
speeds. The standard set $[n-1]=\{1,\ldots,n-1\}$ attains equality.
We call equality \emph{tightness} and throughout set
\[
 R=\{r_1,r_2\}\subset[n-1],\qquad
 V=([n-1]\setminus R)\cup\{u,w\},\qquad n\le u<w.
\]
Zhang~\cite[Questions 8.1--8.2]{Z} asks for a classification and a
polynomial, or $O(n\log n)$, bound on the largest inserted speed.

\begin{theorem}\label{thm:main}
If $n\ge15$ and $\LR(V)\le1/n$, then
\[
 u\le\frac{4n(n-1)^2}{n-4},\qquad w\le2(n-1)u.
\]
As $n\to\infty$, uniformly over such configurations,
$u\le(4+o(1))\,n\log\log n$ and $w=O(n\log n)$. Along an infinite sequence of tight
configurations, $u=(\tfrac12+o(1))\,n\log\log n$. All logarithms are natural.
\end{theorem}

The explicit inequalities give a finite search for each $n\ge15$.
The optimal growth of $w$ and the two-entry classification for all $n$
remain open. We also check $8\le n<15$ directly. Fixing two changes
matters here. If every entry may be changed and common dilations are
allowed, $k[n-1]$ is tight for arbitrarily large $k$. The two-entry
problem at large $n$ has no such scaling freedom.

\newpage
\section{Time intervals exposed by a deletion}
Let $\bad_v=\{t:\dist{vt}\le1/n\}$ on the real line, periodically, and
write $\good_R$ for the open set where all speeds in $[n-1]\setminus R$
have distance strictly above $1/n$. A single subscript $r$ denotes one
deletion. A configuration with $\LR(V)\le1/n$ must cover $\good_R$ by
$\bad_u\cup\bad_w$. Each component of $\bad_v$ has length $2/(nv)$.
For any interval $J$,
\begin{equation}\label{eq:density}
 |\bad_v\cap J|\le 2|J|/n+2/(nv).
\end{equation}
Indeed, split $J$ into periods of length $1/v$ and a remainder. The
remainder is contained in one interval of length $1/v$, of which exactly
$2/(nv)$ is bad. If $\bad_u,\bad_w$ cover a window of length $L$,
\begin{equation}\label{eq:u}
 (n-4)L\le 2/u+2/w\le4/u .
\end{equation}
This is the two-band density method of~\cite[Section 6]{Z}, with the
remainder estimated by a full period.

For a deletion $r\ge2$ write $s=n-r$, and let $F$ be the least integer
at least $s$ coprime to $r$. In particular $s\le F\le s+r-1=n-1$.
The following windows are also contained in the component formulas
of~\cite[Section 2]{Z}. We include the argument for the estimates used here.

\begin{lemma}\label{lem:window}
For any $F\in[s,n-1]$ coprime to $r$, choose $p$ with $Fp\equiv-1\bmod r$. If $2r>n-1$, the interval
$t=p/r+\delta$, $0<\delta<s/(rnF)$, lies in $\good_r$.
If $2r\le n-1$, the same holds for
$1/(2rn)<\delta<s/(rnF)$ whenever the endpoints are in that order.
For the least eligible $F$ the latter interval has length at least $1/(6rn)$.
\end{lemma}
\begin{proof}
Always $\delta<1/(rn)$. If $ip\not\equiv0,-1\bmod r$, the fractional
part of $ip/r$ is between $1/r$ and $1-2/r$, and $0<i\delta<1/r$.
If $ip\equiv-1\bmod r$, then $i\equiv F\bmod r$. Since $F\ge n-r$,
the largest such $i\le n-1$ is $F$. Thus
$i\delta<(n-r)/(rn)=1/r-1/n$. In either case $\dist{it}>1/n$.
Any remaining multiple $i=kr$ has $k\ge2$ and occurs only in the second
case, where $1/n<kr\delta<1/r\le1/2$.

For the length estimate take the least eligible $F$. The distance from any integer to the next coprime
integer, counting distance zero at a unit, is at most $\lfloor r/2\rfloor$.
For odd $r$ use the unit residues $1,(r-1)/2,(r+1)/2,r-1$.
For $4\mid r$ use $1,r/2-1,r/2+1,r-1$; for $r\equiv2\bmod4$, $r\ge6$,
use $1,r/2-2,r/2+2,r-1$. Their cyclic gaps are at most $\lfloor r/2\rfloor+1$;
$r=2$ is immediate. Consequently $F\le s+r/2\le3s/2$ in the second
case. The interval length is $(s/F-1/2)/(rn)\ge1/(6rn)$.
\end{proof}

\newpage
\section{Why one inserted speed cannot fill both gaps}
\begin{lemma}\label{lem:one}
For $n\ge15$ and any integer $v\ge n$, the set $\good_R\setminus\bad_v$
contains a nonempty open interval.
\end{lemma}
\begin{proof}
It suffices to disprove $\good_R\subseteq\bad_v$. If $1\in R$, use the
window $(1/(2n),1/n)$. Its length would force $v\le4$.

Suppose some $r\in R$ satisfies $2\le r\le(n-1)/2$. The long window of
Lemma~\ref{lem:window} would give $v\le12r$. On both sides of $1/r$,
the points $1/r\pm\delta$ also lie in $\good_r$ throughout
\[
 a=\frac1{2rn}<\delta<b=\frac{n-r}{rn(n-1)}.
\]
For nonmultiples use $\dist{i/r}\ge1/r$ and $i\delta<(n-r)/(rn)$;
for surviving multiples use $1/n<i\delta<1/2$. The gap $b-a$ is positive.
If both intervals are covered, their inner endpoints are covered by
closedness. After scaling by $v$, these endpoints are only
$2va=v/(rn)\le12/n<1-2/n$ apart. They must lie in the same closed band
of radius $1/n$ about an integer. By connectedness both full intervals
must stay in that band. Their outer endpoints then give $2vb\le2/n$,
or $v\le r(n-1)/(n-r)<n$.

It remains that $(n-1)/2<r_1<r_2$. At each time $1/r_i$ every retained
speed is a nonmultiple of $r_i$ and has distance at least $1/r_i>1/n$.
Hence coverage forces both $r_1\mid v$ and $r_2\mid v$. Around $1/r_1$
the interval of radius $h=(n-r_1)/(nr_1(n-1))$ is good by the same
Lipschitz estimate. Since its centre maps to an integer under multiplication
by $v$, coverage requires $vh\le1/n$. But with $d=r_2-r_1$ we have
\[
 v\ \ge\ \frac{r_1r_2}{\gcd(r_1,r_2)}
 \ \ge\ \frac{r_1r_2}{d}
 \ >\ \frac{r_1(n-1)}{n-r_1}.
\]
The strict inequality follows from
$r_2(n-r_1)-d(n-1)=r_1(n-r_2)+d>0$. This contradicts $vh\le1/n$.
\end{proof}

\paragraph{Finite cutoff.}
The walls of $\bad_i$ have the form $(nk\pm1)/(ni)$. Distinct walls for
speeds $i,j$ differ by at least $1/(nij)$; walls for the same speed $i$
differ by at least $1/(ni)$. Thus a nonempty component of $\good_R$
has length at least $1/[n(n-1)^2]$. Equation~\eqref{eq:u} bounds $u$.
Lemma~\ref{lem:one} gives a residual window after adding $u$. At most
one of its two different wall speeds can exceed $n-1$, so its length is
at least $1/[n(n-1)u]$, also when both walls have speed $u$.
It must fit in a single band of $\bad_w$. This gives $w\le2(n-1)u$.

\newpage
\section{A double logarithm for the smaller speed}
If either removed speed is one, \eqref{eq:u} on $(1/(2n),1/n)$ gives
$u\le8n/(n-4)$. If a removed speed lies in $[2,(n-1)/2]$, the second
window of Lemma~\ref{lem:window} gives $u\le24rn/(n-4)=O(n)$. Consider
two larger removed speeds. Order their offsets as $s_1<s_2$, put
$r_i=n-s_i$, and let $F_i$ denote the least integer $\ge s_i$ coprime
to $r_i$. Define $K=\min(F_1/s_1,F_2/s_2)$. The first window of
Lemma~\ref{lem:window} and \eqref{eq:u} give
\[
 u\le \min_i\frac{4nr_i F_i}{(n-4)s_i}
 \le \frac{4n^2}{n-4}\,K .
\]
We show $K\le(1+o(1))\log\log n$. Only sequences with $K\to\infty$
require attention, and they have $s_i\ge2$. For each prime $p\le K$ take
the least power $p^a\ge s_i$. Then
$s_i\le p^a<p s_i\le K s_i\le F_i$. By definition of $F_i$ this power
shares a factor with $r_i$, so $p\mid r_i$ for \emph{both} $i$.
Writing $\theta(x)=\sum_{p\le x}\log p$, we obtain
\begin{equation}\label{eq:primorial}
 \exp(\theta(K))\mid (s_2-s_1),\qquad s_2>\exp(\theta(K)).
\end{equation}
Here the first expression is the integer primorial. Moreover every prime
in $[s_2,Ks_2)$ divides $r_2$. By the prime number theorem
$\theta(x)\sim x$, the sum of their logarithms is $(1+o(1))K s_2$ as
$K\to\infty$. Hence $\log n\ge(1+o(1))K s_2$.
Taking another logarithm and using~\eqref{eq:primorial} gives
$\log\log n\ge\theta(K)+\log K+o(1)=(1+o(1))K$. This proves the upper
estimate of Theorem~\ref{thm:main}.

\paragraph{Construction attaining the order.}
Specialize the simultaneous-acceleration construction of Goddyn and
Wong~\cite[Theorem 3.1 and Corollary 3.3]{GW}. For integers $m\ge2$, take
\[
 P=\prod_{p\le2m}p,\qquad s=P+2,\qquad
 Q=\prod_{p\le ms}p,\qquad n=Q+s.
\]
Multiply the two speeds $n-2$ and $n-s=Q$ by $m$. Both inserted
speeds are distinct and at least $n$. The criterion in~\cite{GW} asks
that every integer in $[d,md)$ share a factor with the speed $n-d$ being
multiplied. For $d=2$, use $P\mid Q+P=n-2$; for $d=s$, every integer in
$[s,ms)$ has a prime factor dividing $Q$. Thus the configuration is tight.
For example, $m=2$ gives $n=30038$ and $(u,w)=(60060,60072)$.
Since $\log P\sim2m$, $\log Q\sim ms$ and $s=o(Q)$, we get
$\log\log n\sim 2m$ and $u=mQ\sim mn$. The construction is a quantitative
specialization of the earlier scheme; the order bound applies also to
configurations with no assumed multiplicative pairing.

\newpage
\section{A single logarithm for the larger speed}
Suppose, seeking a contradiction, there is a sequence with $n\to\infty$
and $w/(n\log n)\to\infty$. The previous section gives $u=O(n\log\log n)$.
Every open interval in $\good_R$ disjoint from $\bad_u$ must have length
at most $2/(nw)$, since it fits in one connected component of $\bad_w$.
We apply this observation at several numerators $p$ for the same deletion.

\paragraph{A deletion in the lower half.}
The case $r=1$ is excluded for large $n$ by the bound in Section~4.
If $2\le r\le(n-1)/2$ along a subsequence, Section~4 gives
$n\le u\le24rn/(n-4)$. In particular $r$ and $s=n-r$ are both of order
$n$, and $u/r$ is bounded. Use all primes $F$ in
$(s,s+\min(r/2,s/10)]$. There are unboundedly many by the prime number
theorem, as the interval length is bounded below by a fixed multiple
of $s$. All exceed $r$, are distinct modulo $r$, and are at most $n-1$.
Each gives a window from Lemma~\ref{lem:window} of length at least
$(10/11-1/2)/(rn)>2/(nw)$. Thus each window meets $\bad_u$.
At its centre of reference $p/r$ the Lipschitz bound consequently gives
\[
 \dist{up/r}\le 1/n+u/(rn).
\]
The least absolute residue $j$ of $up$ modulo $r$ therefore satisfies
$|j|\le (r+u)/n=O(1)$ and $jF\equiv-u\bmod r$. If $r\nmid u$ then
$j\ne0$. For each fixed nonzero $j$ at most $\gcd(|j|,r)\le|j|$
residue classes $F$ solve this congruence, contradicting the number
of available primes. Hence $u=mr$ with $3\le m=O(1)$, since $2r<n$.
But throughout each of the same windows, $ut$ has distance
$mr\delta$ to the nearest integer: $1/n<m/(2n)<mr\delta<m/n<1/2$
for large $n$. This leaves a whole window for $w$, a contradiction.

\paragraph{Divisibility in the upper half.}
We may now suppose both deleted speeds exceed $(n-1)/2$. Fix either
$r$, put $s=n-r$, and set $B=\min(n-1,20s\log n)$. For large $n$,
$(3B/5,B]$ lies above $s$, and its primes have total log weight
$(2/5+o(1))B$ by the prime number theorem. Since $B\ge20\log n$
eventually, this exceeds $\log n+2\log B$ (use $\log B=o(B)$).
The primes dividing $r$ contribute at most $\log r$. Hence at least
three distinct primes $F$ in this interval are coprime to $r$. They lie
in $[s,n-1]$, a set of $r$ successive integers.

For each such $F$ choose $Fp\equiv-1\bmod r$. The one-sided window at
$p/r$ has length $s/(rnF)\ge1/(20rn\log n)$. If $\dist{up/r}\ge2/r$,
the initial portion of length
\[
 \min\{s/(rnF),\,1/(nu)\}
\]
is disjoint from $\bad_u$: the distance at the centre exceeds $1/n$
by at least $1/n$. Both entries of the minimum exceed $2/(nw)$ for large
$n$. This contradicts coverage. Thus the least absolute residue of $up$
is at most one for every chosen $F$. Unless $r\mid u$, this forces
$F\equiv\pm u\bmod r$ for three distinct residue classes. We conclude
that both deleted speeds divide $u$.

\newpage
\paragraph{The two forced divisibilities are incompatible with coverage.}
Write $r_1<r_2$ and $d=r_2-r_1$. Then
\[
 u\ge\frac{r_1r_2}{\gcd(r_1,r_2)}\ge\frac{r_1r_2}{d}.
\]
Consequently $s=n-r_1>d\gg n/\log\log n$, and $s\le n/2$.
The prime number theorem supplies a prime $F$ in $(s,6s/5]$
which does not divide $r_1$: the log weight of those primes is
$(1/5+o(1))s>\log n$. In particular $F\le n-1$ for large $n$.
Lemma~\ref{lem:window} at the corresponding $p/r_1$ yields a free
interval containing
\[
 t=p/r_1+\delta,\qquad
 \frac{3}{5r_1n}<\delta<\frac{3}{4r_1n},
\]
since $s/F\ge5/6$. Now $u=mr_1$ with $2\le m=O(\log\log n)$.
On this interval $3m/(5n)<u\delta<3m/(4n)$,
so $1/n<u\delta<1/2$ for large $n$. Thus the interval is disjoint
from $\bad_u$, and has length $3/(20r_1n)>2/(nw)$.
This final contradiction proves $w=O(n\log n)$.

\paragraph{Where the two logarithmic scales arise.}
For the smaller speed, both deletions must have long gaps to the next
coprime integer. The resulting common divisor forces the offsets far
apart. For the larger speed, three coprime choices in a short enough
range give three distinct rational tests. A speed exceeding the claimed
order would leave the smaller speed responsible for all three at each
deletion, forcing two divisibilities that the last interval rules out.
The construction in Section~4 attains $n\log\log n$ for both speeds.
It leaves open whether the larger speed can grow as fast as $n\log n$.

\newpage
\section{Exact census with unbounded inserted speeds}
For each pair of deleted speeds we work on the half circle by symmetry,
in the coordinate $x=nt$. The walls become $(nk\pm1)/v$. Sweep the
original speeds and record the open cells on which at most two are bad.
After deleting $\{r,s\}$, retain cells whose bad speeds form a subset of
$\{r,s\}$. Dirichlet's pigeonhole argument for $0,t,\ldots,(n-1)t$
ensures at least one speed is bad. Merge adjacent free cells; the covering
bands are closed, so checking the interiors suffices.

If $L$ is the largest free length in the $x$ coordinate, test integers
$n\le u\le\lfloor 4n/((n-4)L)\rfloor$ by~\eqref{eq:u}. Remove the
closed bands of $u$. If the longest remaining interval has length $D$,
test $u<w\le\lfloor 2/D\rfloor$. The program explicitly checks that the
residual is nonempty, including at $8\le n<15$. For each interval $(a,b)$,
one closed band of $w$ suffices exactly when an integer $j$ satisfies
$wb-1\le nj\le wa+1$. All decisions use integer arithmetic. A strict
witness at a wall would persist in an adjacent open interval.

\begin{center}
\begin{tabular}{rrrrr}
\toprule $n$&\multicolumn{2}{c}{Deleted}&\multicolumn{2}{c}{Inserted}\\
\midrule
8&2&3&11&13\\74&70&72&140&144\\284&280&282&560&564\\
\bottomrule
\end{tabular}
\end{center}
These are all configurations with $\LR(V)\le1/n$ in $8\le n\le300$.
In each surviving case $t=1/n$ witnesses equality. The scan treats
$4\,455\,065$ deletion pairs. A separate implementation with rational
interval intersections checks all sizes $8$ through $18$ from scratch
and directly verifies all three survivors. Zhang~\cite[Theorem 6.10]{Z}
previously covered $8\le n\le140$ with $w\le8n$; the third row was
already identified in~\cite[Example 5.2]{Z}. The additional information
here is the exhaustive range in $n$ and the absence of a prescribed speed
ceiling. A finite census does not establish the classification for all $n$.

\paragraph{Reproducibility and status.}
The supplement includes the C++17 enumeration, per-size counts, and
the Python rational cross-check. This draft and the software were
prepared using OpenAI Codex. Independent mathematical review and
references to earlier equivalent bounds are welcome.

\begin{thebibliography}{9}\small\raggedright
\bibitem{GW} L. Goddyn and E. B. Wong, Tight instances of the lonely runner,
\textit{Integers} \textbf{6} (2006), A38.
\href{https://www.sfu.ca/~goddyn/Papers/063tight_lonely_runner.pdf}{Author's copy}.
\bibitem{Z} Y. Zhang, Tight instances of the Lonely Runner Conjecture:
one-entry modifications and the growth of speeds, arXiv:2608.13599v3 (2026).
\href{https://arxiv.org/html/2608.13599v3}{Version 3, Sections 2, 5, 6 and 8}.
\end{thebibliography}
\end{document}
